% DF-001, CF-007/012, P09 del inventario REV05.
% Residencia: APP toroidal y transporte de cursor; ejemplo tras el cociclo.
% Propietarios: PAQUETE_ARTICULO_I_INTEGRACION_ESTADO_CAMPOS_20260919/
% deltas/integracion/APPRouteWinding.lean (íntegro),
% lean/biblioteca/TPKTransport.lean:19--143, 155--229,
% lean/biblioteca/TPKFiniteCursor.lean:17--146.
% Estatuto: RESULTADO_RECUPERADO; desarrollo del ejemplo finito efectivo,
% sin identificar el factor observable con todo el estado enriquecido.

\section[Spatial return and winding of the cursor]{Spatial return and conservation of winding in the calendar of twenty-seven transitions}
\label{rc27:sec:cursor}

The two periodic APP coordinates arise from two transported integer
coordinates. The residual reading identifies positions that
differ by multiples of nine; the integer quotients preserve the
turns of that transport. The TPK calendar provides a complete
example in which the residual position returns and the lifted cursor
preserves a nonzero horizontal displacement.

\subsection{Lifted coordinates and boundary cocycle}

For $z\in\mathbb Z$, let
\[
\rho_9(z)=1+((z-1)\bmod9),\qquad
q_9(z)=\left\lfloor\frac{z-1}{9}\right\rfloor.
\]
Euclidean division, also for negative coordinates, gives
\begin{equation}
z=\rho_9(z)+9q_9(z),\qquad1\le\rho_9(z)\le9.
\label{rc27:eq:division}
\end{equation}
The transport carry of an increment $a\in\mathbb Z$ is
\[
c(z,a)=q_9(z+a)-q_9(z).
\]
From \eqref{rc27:eq:division} we deduce
\begin{align}
\rho_9(z+a)+9c(z,a)&=\rho_9(z)+a,\\
c(z,a+b)&=c(z,a)+c(z+a,b).
\label{rc27:eq:cociclo}
\end{align}
The second identity follows by canceling $q_9(z+a)$ on the right-hand
side. The second increment is evaluated from the endpoint of the first;
that dependence preserves the boundary crossing information.

Let $D=\{N,E,S,O\}$, with displacements
\[
\delta(N)=(-1,0),\quad\delta(E)=(0,1),\quad
\delta(S)=(1,0),\quad\delta(O)=(0,-1).
\]
For a direction word $w=d_1\cdots d_n$, define
$\Delta(w)=\sum_j\delta(d_j)$ and, from $o=(o_1,o_2)$,
\[
\operatorname{Ext}(o,w)=o+\Delta(w),\qquad
\nu(o,w)=\bigl(c(o_1,\Delta_1(w)),c(o_2,\Delta_2(w))\bigr).
\]
The vector $\nu$ records both windings.

\begin{proposition}[Composition and reversal of turns]
For two concatenable routes $u,v$ one has
\[
\nu(o,uv)=\nu(o,u)+\nu(\operatorname{Ext}(o,u),v).
\]
If $w^{-1}$ reverses the order and replaces each direction with its
opposite, then
\[
\operatorname{Ext}(\operatorname{Ext}(o,w),w^{-1})=o,
\qquad
\nu(\operatorname{Ext}(o,w),w^{-1})=-\nu(o,w).
\]
For a spatially closed route in APP,
\[
\nu(o,w)=\Delta(w)/9\in\mathbb Z^2.
\]
\end{proposition}
\begin{proof}
The sum of displacements is additive under concatenation. Applying
\eqref{rc27:eq:cociclo} in each coordinate proves the first
identity. Reversal gives $\Delta(w^{-1})=-\Delta(w)$ and recovers
the initial endpoint. Subtracting the quotients in reverse order changes
the sign of $\nu$. Finally, spatial closure is equivalent to
$\rho_9(o_i+\Delta_i)=\rho_9(o_i)$ for both coordinates; the first
equality of \eqref{rc27:eq:cociclo} then gives $9\nu_i=\Delta_i$.
\end{proof}

The cycles of nine southward moves and nine eastward moves,
from $(1,1)$, have turns $(1,0)$ and $(0,1)$, respectively.
They are two independent cycles of the periodic APP surface.

\subsection{Update of the two cursors}

A lifted observable cursor is $c=(x,y,d)\in\mathbb Z^2\times D$.
The phase $p$ belongs to $\mathbb Z/108\mathbb Z$ and is represented
by $0,\ldots,107$. In the additive sheet, displacement is activated
when $p\equiv0\pmod3$; in the multiplicative sheet, when
$p\equiv1\pmod3$. Define
\[
a_+(p)=\mathbf1_{\{p\equiv0\ (3)\}},\qquad
a_\times(p)=\mathbf1_{\{p\equiv1\ (3)\}},\qquad
\varepsilon(p)=\mathbf1_{\{p+1\equiv0\ (27)\}}.
\]
The observable operator acts through
\begin{equation}
c'_s=\bigl(x_s+a_s(p)\delta_1(d_s),\,
           y_s+a_s(p)\delta_2(d_s),\,
           \operatorname{op}^{\varepsilon(p)}d_s\bigr),
\qquad p'=p+1\pmod{108},
\label{rc27:eq:actualizacion}
\end{equation}
where $s\in\{+,\times\}$ and $\operatorname{op}$ interchanges
$N$ with $S$ and $E$ with $O$.
Displacement uses the preceding direction and then the
half-cycle transition updates the orientation.

The inverse first recovers $p=p'-1\pmod{108}$, restores
$d_s=\operatorname{op}^{\varepsilon(p)}d'_s$ and subtracts
$a_s(p)\delta(d_s)$ from the coordinates. Since
$\operatorname{op}^2=I$, both compositions are the identity.

\begin{theorem}[Effective spatial return in the first half-crown]
\label{rc27:thm:ejemplo}
Start from
$c_+(0)=c_\times(0)=(1,1,E)$ and $p(0)=0$.
After twenty-seven transitions of
\eqref{rc27:eq:actualizacion}, one obtains
\[
c_+(27)=c_\times(27)=(1,10,O),\qquad p(27)=27.
\]
Both residual positions return to $(1,1)$, while the
horizontal quotient increases by one and the vertical quotient remains
unchanged.
\end{theorem}
\begin{proof}
For $0\le n\le27$, the number of additive activations before
transition $n$ is $\lfloor(n+2)/3\rfloor$; the number of
multiplicative activations is $\lfloor(n+1)/3\rfloor$. Before
the phase transition from $26$ to $27$ there is no direction reversal.
By induction on $n$, the coordinates are
\[
x_+(n)=x_\times(n)=1,\qquad
y_+(n)=1+\left\lfloor\frac{n+2}{3}\right\rfloor,\qquad
y_\times(n)=1+\left\lfloor\frac{n+1}{3}\right\rfloor.
\]
The directions are $E$ for $0\le n<27$ and $O$ for $n=27$.
The decomposition in each group of three transitions is
\[
\begin{array}{c|c|c|c}
n&y_+(n)&y_\times(n)&\text{direction}\\\hline
3k&1+k&1+k&E\\
3k+1&2+k&1+k&E\\
3k+2&2+k&2+k&E
\end{array}
\quad(0\le k\le8),
\]
and the final endpoint is $(y_+,y_\times,d)=(10,10,O)$.
There are nine activations in each sheet. The final transition has preceding
phase $26$, of class two modulo three; it therefore reverses the
directions and maintains the coordinates that have already reached ten.

Finally, $\rho_9(10)=1$, $q_9(10)=1$,
$\rho_9(1)=1$ and $q_9(1)=0$. The residual position coincides with
the initial one, and the spatial winding is $(0,1)$.
\end{proof}

The return in this example is spatial. The recorded direction and phase
are $(O,27)$ at the end and $(E,0)$ at the start. The example
expressly preserves the three data: residual position, orientation and
phase; their equality or difference is established for each one.

\subsection{Cursor recovery and compatibility at every depth}

\begin{proposition}[Recovery of the lifted cursor]
The residual position, the direction and the two integer quotients
uniquely determine the cursor $c=(x,y,d)$.
\end{proposition}
\begin{proof}
The two identities \eqref{rc27:eq:division} recover $x$ and $y$
and the direction is preserved in the observable projection. If these
five data coincide for two cursors, their three
components coincide and the cursors are therefore equal.
\end{proof}

Let $P$ be the projection of the two cursors onto their positions modulo
nine, preserving their directions and the phase. The finite
update $T_{\rm fin}$ uses the same activation and
orientation indicators, with residual addition in the coordinates. If $T$ is
\eqref{rc27:eq:actualizacion}, then
\[
PT=T_{\rm fin}P,\qquad PT^n=T_{\rm fin}^{\,n}P
\quad(n\ge0).
\]
The first equality results from compatibility of reduction
modulo nine with addition and preservation of direction and phase;
the second is proved by induction on $n$.

The ordered archive of observable states can be prolonged through
$(q,\mathcal A)\mapsto(Tq,\mathcal A\Vert q)$.
After $n$ transitions it preserves its initial prefix and its length
increases by $n$. This operation keeps the current state,
the winding and the register of its preceding states separate. In the
complete enriched construction, the observable factor is composed
with the other cylinder, incidence and memory readings through their
declared projections.
